\chapter{Problem Statement, Analysis, and Motivation}
\label{ch:motivation}
This chapter develops the motivation for addressing temporary data gaps, explains why the problem becomes particularly severe in SoTS, and outlines why DTs provide an advantageous foundation for mitigating this challenge. The chapter also identifies the limitations of existing event-driven processing approaches, drawing on insights from the systematic literature review in \chref{ch:literature-review}, and defines the scope of the architectural solution presented in \chref{ch:architecture}.

\section{Temporary Data Gaps in Event-Driven Computation}

A temporary data gap occurs when an observation expected at a known interval does not arrive in time for the computation that depends on it. These gaps interrupt the flow of information through an event-driven pipeline, forcing downstream operators to reason with stale or incomplete context. Pattern-based reasoning engines fail to match temporal patterns, and systems dependent on regularly synchronized views lose alignment with the physical processes they monitor. Any streaming system whose computation assumes roughly periodic updates is therefore sensitive to gaps. In loosely coupled settings the impact may remain localized, but in tightly integrated systems even brief lapses can propagate through the pipeline, degrading state, pattern recognition, and system responsiveness.

\section{Why SoTS Are Especially Vulnerable}

SoTS embody the kind of tightly connected, interdependent environment in which temporary data gaps are both common and disruptive. As shown in the literature review, existing SoTS implementations frequently rely on continuous cross-twin communication and shared situational awareness, yet they provide limited mechanisms to manage the consequences of temporary disconnections or configuration changes. The review further highlights that runtime reconfiguration, intermittent connectivity, and constituent mobility are recurring characteristics of SoTS, but existing work rarely addresses the resulting instability. At the same time, SoTS depend heavily on coordinated reasoning across multiple DTs, meaning that a data gap affecting one twin can readily propagate to others. This combination of high interdependence and limited runtime support for dealing with data discontinuities makes temporary data gaps a structural challenge in SoTS. Addressing these gaps is therefore essential for enabling reliable operation in SoTS, particularly as they evolve toward more dynamic, adaptive, and distributed forms.

\section{Problem Statement}

The preceding discussion leads to a central problem:

\begin{quote}
\emph{Event-driven systems, and SoTS in particular, lack mechanisms for allowing event-driven reasoning to operate during intervals in which expected observations temporarily fail to arrive.}
\end{quote}

What is missing is an architectural means to preserve the continuity of event-driven computation during such gaps. Specifically, there is no principled way to represent what should have been observed, how confident the system is in that representation, or how downstream reasoning should proceed in the absence of timely data. Without such a mechanism, SoTS remain vulnerable to momentary lapses in data availability, even when they possess the modelling and simulation capabilities that could, in principle, mitigate these effects.

\section{Limitations of Existing Event-Driven Approaches}

Although temporary data gaps arise across many classes of event-driven computation, existing systems offer limited support for representing or reasoning through intervals in which data is unavailable. Most event-driven architectures assume that inputs arrive predictably, and downstream operators are not designed to operate reliably when those expected updates fail to appear.

These limitations are especially evident in Complex Event Processing (CEP). As summarized by \citet{alevizos2017probabilistic}, probabilistic CEP frameworks address uncertainty only in events that \emph{arrive}. They accommodate noisy attributes, uncertain timestamps, or rule-level ambiguity but assume uninterrupted event flows. Even approaches that infer hypothetical subevents operate within recognition logic rather than modeling the absence of expected observations. None provide semantics for detecting, reconstructing, or propagating uncertainty over missing events. Other CEP extensions reinforce this gap: Bayesian-network models \cite{wasserkrug2010efficient, wasserkrug2012model}, interval-based timestamp reasoning \cite{zhang2013recognizing}, and lineage- or NFA-based probabilistic engines \cite{shen2008probabilistic, kawashima2010complex, chuanfei2010complex, wang2013complex} all require the presence of input events to progress. Frameworks such as CEP2U \cite{cugola2015introducing} and uncertainty-aware extensions to Esper and Flink \cite{moreno2019managing} enrich reasoning about noisy or partially observed events but remain inactive when the stream temporarily disappears.

A small number of IoT-oriented studies attempt to mitigate missing data through external fabrication or interpolation. \citet{gkoulis2025assessing} evaluate naïve and trend-aware smoothing techniques, and \citet{gkoulis2024hybrid} introduce a quality-aware CEP platform that injects synthetic events. While these approaches acknowledge the practical consequences of missing data, reconstruction is performed outside the CEP semantics: the event processor cannot distinguish synthetic from genuine inputs, nor can it propagate uncertainty about the fabricated values.

More broadly, the literature review shows that reliability in SoTS is often discussed but rarely formalized or experimentally validated, and almost no existing work addresses reliability under intermittent data availability. Current event-driven systems, including those embedded in DT or SoTS architectures, offer no integrated mechanism for reconstructing expected-but-missing observations or for sustaining reasoning when the input stream falls silent. This absence motivates the architectural method developed in this thesis.

\section{Why Digital Twins Are Uniquely Positioned to Mitigate Data Gaps}

Meeting this problem requires a capability that existing event-driven systems lack: maintaining a usable view of system state when observations momentarily disappear. Digital twins provide this through their modelling, forecasting, and state-estimation capacities, which allow them to infer short-term behaviour during brief interruptions and supply estimates that remain temporally consistent with the underlying process. However, as identified in the literature review, current DT implementations generally use predictive models for synchronization, simulation, or control rather than for sustaining event-driven computation. Existing platforms do not provide a formal mechanism for reinserting forecasts into event streams or integrating them with high-level reasoning engines such as CEP. As a result, the predictive abilities of DTs remain internal utilities rather than components that actively maintain continuity in event-driven processing. The architectural contribution of this thesis is to establish that role: a method for integrating DT-driven reconstruction with event-driven reasoning, making prediction a mechanism for maintaining reliable event processing.

\section{Scope of the Architecture}

The architectural solution proposed in this thesis focuses specifically on short-duration disruptions in event availability. The solution does not target long-term outages, permanent sensor failures, structurally missing data, or events with corrupted or invalid payloads. In such cases, reconstruction is neither feasible nor aligned with the goal of preserving continuity during brief periods of silence. The architecture also acknowledges that real-world event streams may include late, out-of-order, or backlogged events, but these conditions do not trigger rollback or retroactive revision of previously issued outputs. Delayed events arriving after a reconstruction interval are forwarded immediately, without modifying earlier computations. Out-of-order events are inserted according to their source timestamps, but past outputs remain fixed to avoid destabilizing downstream reasoning. When a generator reconnects and emits buffered events, they are passed through intact. Reconstructions remain part of the logical history. These late arrivals may be used to recalibrate predictive models but do not alter earlier decisions.


\section{Summary}

This chapter has shown that temporary data gaps pose a fundamental challenge for event-driven computation and that their effects are amplified in SoTS, where continuous cross-twin dataflows are assumed. The literature review indicates that current SoTS offer little support for handling short-term disruptions caused by runtime reconfiguration or intermittent connectivity. Existing event-driven systems, including uncertainty-aware CEP, advance their reasoning only when new events arrive and therefore cannot detect, represent, or compensate for intervals in which expected observations fail to appear. Digital twins possess modelling and forecasting capabilities that could supply estimates during these interruptions, but current platforms keep these capabilities internal and do not integrate them into event-driven reasoning pipelines. As a result, SoTS lack a mechanism for allowing CEP to operate when data is temporarily unavailable. This motivates the architectural approach introduced in the next chapter (\chref{ch:architecture}), which incorporates DT-driven reconstruction into event-driven reasoning to handle temporary data gaps explicitly.

